\rSec2[expected.ref]{Partial specialization of expected for reference types}

\rSec3[expected.ref.general]{General}

\begin{codeblock}
template<class T, class E>
class @\libglobal{expected}@<T&, E> {
private:
  using @\exposidnc{error-value-type}@ = remove_cv_t<remove_reference_t<E>>; // exposition only

public:
  using @\libmember{value_type}{expected}@ = T&;
  using @\libmember{error_type}{expected}@ = E;
  using @\libmember{unexpected_type}{expected}@ = unexpected<E>;

  template<class U> using rebind = expected<U, error_type>;

  // -------------------------------------------------------------------------
  // Constructors
  // -------------------------------------------------------------------------

  expected() = BEMAN_EXPECTED_DELETE_MSG(
      "expected<T&,E>: no default constructor; T& cannot be null");

  // Copy constructor (trivial path). Unconstrained; see the primary
  // template's copy constructor for why.
  constexpr expected(const expected&) = default;

  // Copy constructor (non-trivial path)
  constexpr expected(const expected& rhs) noexcept(is_nothrow_copy_constructible_v<E>)
    requires(is_copy_constructible_v<E> && !is_trivially_copy_constructible_v<E>);

  // Move constructor (trivial path). Unconstrained; no explicit noexcept.
  constexpr expected(expected&&) = default;

  // Move constructor (non-trivial path)
  constexpr expected(expected&& rhs) noexcept(is_nothrow_move_constructible_v<E>)
    requires(is_move_constructible_v<E> && !is_trivially_move_constructible_v<E>);

  // Deleted: no in-place value constructor — T& cannot be constructed in-place
  template<class... Args>
  constexpr expected(in_place_t, Args&&...) = BEMAN_EXPECTED_DELETE_MSG(
      "expected<T&,E>: no in-place value constructor; T& cannot be constructed "
      "in-place — pass a U convertible to T&");

  // Value constructor — takes U that can bind to T&
  template<class U = T>
    requires(!is_same_v<remove_cvref_t<U>, in_place_t> &&
             !is_same_v<remove_cvref_t<U>, expected<T&, E>> &&
             !@\exposidnc{is-unexpected-specialization}@<remove_cvref_t<U>>::value &&
             is_constructible_v<T&, U> && !reference_constructs_from_temporary_v<T&, U>)
  constexpr explicit(!is_convertible_v<U, T&>)
      expected(U&& u) noexcept(is_nothrow_constructible_v<T&, U>);

  // Deleted: binding a temporary to T& creates a dangling reference
  template<class U>
    requires(reference_constructs_from_temporary_v<T&, U>)
  constexpr expected(U&&) =
      BEMAN_EXPECTED_DELETE_MSG("expected<T&,E>: argument would bind a temporary that "
                                "dangles; pass an lvalue reference");

  // Converting constructor from expected<U&, G> (copy) — value-E path
  template<class U, class G>
    requires(!is_reference_v<E> && is_constructible_v<T&, U&> &&
             is_constructible_v<E, const G&> &&
             !reference_constructs_from_temporary_v<T&, U&>)
  constexpr explicit(!is_convertible_v<U&, T&> || !is_convertible_v<const G&, E>)
      expected(const expected<U&, G>& rhs);

  // Converting constructor from expected<U&, G> (move) — value-E path
  template<class U, class G>
    requires(!is_reference_v<E> && is_constructible_v<T&, U&> &&
             is_constructible_v<E, G> && !reference_constructs_from_temporary_v<T&, U&>)
  constexpr explicit(!is_convertible_v<U&, T&> || !is_convertible_v<G, E>)
      expected(expected<U&, G>&& rhs);

  // Converting constructor from expected<U&, G&> (copy/move) — reference-E path: only
  // accepts sources whose error type G is itself a reference convertible to E.
  template<class U, class G>
    requires(is_reference_v<E> && is_reference_v<G> && is_constructible_v<T&, U&> &&
             is_convertible_v<G, E> && !reference_constructs_from_temporary_v<T&, U&>)
  constexpr explicit(!is_convertible_v<U&, T&> || !is_convertible_v<G, E>)
      expected(const expected<U&, G>& rhs);

  template<class U, class G>
    requires(is_reference_v<E> && is_reference_v<G> && is_constructible_v<T&, U&> &&
             is_convertible_v<G, E> && !reference_constructs_from_temporary_v<T&, U&>)
  constexpr explicit(!is_convertible_v<U&, T&> || !is_convertible_v<G, E>)
      expected(expected<U&, G>&& rhs);

  // Constructor from unexpected<G> const& / && — value-E path
  template<class G>
    requires(!is_reference_v<E> && is_constructible_v<E, const G&>)
  constexpr explicit(!is_convertible_v<const G&, E>) expected(const unexpected<G>& e);

  template<class G>
    requires(!is_reference_v<E> && is_constructible_v<E, G>)
  constexpr explicit(!is_convertible_v<G, E>) expected(unexpected<G>&& e);

  // Constructor from unexpected<G> — reference-E path. Allowed only when G is itself a
  // reference, i.e. the source unexpected holds a reference to an external object, so
  // binding E& to e.error() cannot dangle regardless of the source's value category. No
  // const_cast is needed.
  template<class G>
    requires(is_reference_v<E> && is_reference_v<G> && is_constructible_v<E, G> &&
             !reference_constructs_from_temporary_v<E, G>)
  constexpr explicit(!is_convertible_v<G, E>) expected(const unexpected<G>& e) noexcept;

  template<class G>
    requires(is_reference_v<E> && is_reference_v<G> && is_constructible_v<E, G> &&
             !reference_constructs_from_temporary_v<E, G>)
  constexpr explicit(!is_convertible_v<G, E>) expected(unexpected<G>&& e) noexcept;

  // Deleted for reference E with value G: the referent lives inside the unexpected<G>
  // object, so binding E& to it would dangle once a temporary source is destroyed. Use
  // (unexpect, lvalue), or an unexpected<E&> holding an external object, instead.
  template<class G>
    requires(is_reference_v<E> && !is_reference_v<G>)
  constexpr expected(const unexpected<G>&) = BEMAN_EXPECTED_DELETE_MSG(
      "expected<T&,E&>: cannot construct from unexpected<value>; the value would "
      "dangle — use unexpected<E&>");

  template<class G>
    requires(is_reference_v<E> && !is_reference_v<G>)
  constexpr expected(unexpected<G>&&) = BEMAN_EXPECTED_DELETE_MSG(
      "expected<T&,E&>: cannot construct from unexpected<value>; the value would "
      "dangle — use unexpected<E&>");

  // In-place constructor for error
  template<class... Args>
    requires(is_constructible_v<E, Args...> && !@\exposidnc{unexpect-dangles-v}@<E, Args...>)
  constexpr explicit expected(unexpect_t, Args&&... args);

  // Deleted: single argument would bind E& to a temporary — dangling prevention
  template<class... Args>
    requires(@\exposidnc{unexpect-dangles-v}@<E, Args...>)
  constexpr expected(unexpect_t, Args&&...) =
      BEMAN_EXPECTED_DELETE_MSG("expected<T&,E&>: unexpect argument would bind a "
                                "temporary that dangles; pass an lvalue reference");

  // Deleted catch-all: reference E, argument neither constructible nor a dangling case
  template<class... Args>
    requires(is_reference_v<E> && !is_constructible_v<E, Args...> &&
             !@\exposidnc{unexpect-dangles-v}@<E, Args...>)
  constexpr expected(unexpect_t, Args&&...) = BEMAN_EXPECTED_DELETE_MSG(
      "expected<T&,E&>: no viable conversion from the given argument(s) to E&");

  // In-place constructor for error with initializer_list
  template<class U, class... Args>
    requires(!is_reference_v<E> && is_constructible_v<E, initializer_list<U>&, Args...>)
  constexpr explicit expected(unexpect_t, initializer_list<U> il, Args&&... args);

  template<class U, class... Args>
    requires is_reference_v<E>
  constexpr expected(unexpect_t, initializer_list<U>, Args&&...) =
      BEMAN_EXPECTED_DELETE_MSG("expected<T&,E&>: initializer-list error construction "
                                "cannot bind a reference; pass an lvalue reference");

  // -------------------------------------------------------------------------
  // Destructor
  // -------------------------------------------------------------------------

  constexpr ~expected()
    requires is_trivially_destructible_v<E>
  = default;

  constexpr ~expected()
    requires(!is_trivially_destructible_v<E>);

  // -------------------------------------------------------------------------
  // Assignment (rebind semantics)
  // -------------------------------------------------------------------------

  // Copy assignment (trivial path)
  constexpr expected& operator=(const expected&)
    requires(is_trivially_copy_constructible_v<E> &&
             is_trivially_copy_assignable_v<E> && is_trivially_destructible_v<E>)
  = default;

  // Copy assignment (non-trivial path)
  constexpr expected& operator=(const expected& rhs) noexcept(
      is_nothrow_copy_constructible_v<E> && is_nothrow_copy_assignable_v<E>)
    requires((is_reference_v<E> ||
              (is_copy_constructible_v<E> && is_copy_assignable_v<E>)) &&
             !(is_trivially_copy_constructible_v<E> &&
               is_trivially_copy_assignable_v<E> && is_trivially_destructible_v<E>));

  // Move assignment (trivial path)
  constexpr expected& operator=(expected&&) noexcept
    requires(is_trivially_move_constructible_v<E> &&
             is_trivially_move_assignable_v<E> && is_trivially_destructible_v<E>)
  = default;

  // Move assignment (non-trivial path)
  constexpr expected& operator=(expected&& rhs) noexcept(
      is_nothrow_move_constructible_v<E> && is_nothrow_move_assignable_v<E>)
    requires((is_reference_v<E> ||
              (is_move_constructible_v<E> && is_move_assignable_v<E>)) &&
             !(is_trivially_move_constructible_v<E> &&
               is_trivially_move_assignable_v<E> && is_trivially_destructible_v<E>));

  // Rebind reference from lvalue
  template<class U = T>
    requires(!is_same_v<remove_cvref_t<U>, expected<T&, E>> &&
             !@\exposidnc{is-unexpected-specialization}@<remove_cvref_t<U>>::value &&
             is_constructible_v<T&, U> && !reference_constructs_from_temporary_v<T&, U>)
  constexpr expected& operator=(U&& u);

  // Assignment from unexpected<G> — value-E path
  template<class G>
    requires(!is_reference_v<E> && is_constructible_v<E, const G&> &&
             is_assignable_v<E&, const G&>)
  constexpr expected& operator=(const unexpected<G>& e);

  template<class G>
    requires(!is_reference_v<E> && is_constructible_v<E, G> && is_assignable_v<E&, G>)
  constexpr expected& operator=(unexpected<G>&& e);

  // Rebinding assignment for reference E from reference G — binds unex_ to the external
  // referent.
  template<class G>
    requires(is_reference_v<E> && is_reference_v<G> && is_constructible_v<E, G> &&
             !reference_constructs_from_temporary_v<E, G>)
  constexpr expected& operator=(const unexpected<G>& e);

  template<class G>
    requires(is_reference_v<E> && is_reference_v<G> && is_constructible_v<E, G> &&
             !reference_constructs_from_temporary_v<E, G>)
  constexpr expected& operator=(unexpected<G>&& e);

  // Deleted for reference E with value G: would rebind E& to unexpected<G>'s temporary
  // storage.
  template<class G>
    requires(is_reference_v<E> && !is_reference_v<G>)
  constexpr expected& operator=(const unexpected<G>&) = BEMAN_EXPECTED_DELETE_MSG(
      "expected<T&,E&>: cannot assign from unexpected<value>; the value would dangle — "
      "use unexpected<E&>");

  template<class G>
    requires(is_reference_v<E> && !is_reference_v<G>)
  constexpr expected& operator=(unexpected<G>&&) = BEMAN_EXPECTED_DELETE_MSG(
      "expected<T&,E&>: cannot assign from unexpected<value>; the value would dangle — "
      "use unexpected<E&>");

  // emplace — rebind the reference
  template<class U = T>
    requires(is_constructible_v<T&, U> && !reference_constructs_from_temporary_v<T&, U>)
  constexpr T& emplace(U&& u) noexcept(is_nothrow_constructible_v<T&, U>);

  // -------------------------------------------------------------------------
  // Swap
  // -------------------------------------------------------------------------

  constexpr void swap(expected& rhs) noexcept(is_nothrow_move_constructible_v<E> &&
                                              (is_reference_v<E> ||
                                               is_nothrow_swappable_v<E>))
    requires((is_reference_v<E> || is_swappable_v<E>) && is_move_constructible_v<E>);

  friend constexpr void swap(expected& x, expected& y) noexcept(noexcept(x.swap(y)))
    requires((is_reference_v<E> || is_swappable_v<E>) && is_move_constructible_v<E>);

  // -------------------------------------------------------------------------
  // Observers
  // -------------------------------------------------------------------------

  constexpr T* operator->() const noexcept;
  constexpr T& operator*() const noexcept;

  constexpr explicit operator bool() const noexcept;
  constexpr bool has_value() const noexcept;

  constexpr T& value() const&;
  constexpr T& value() &&;

  // error() — shallow const: always returns E& regardless of const on expected
  constexpr const E& error() const& noexcept;
  constexpr E& error() & noexcept;
  constexpr const E&& error() const&& noexcept;
  constexpr E&& error() && noexcept;

  template<class U = remove_cv_t<T>>
    requires(is_object_v<T> && !is_array_v<T>)
  constexpr remove_cv_t<T> value_or(U&& def) const;

  // Constraints spell error_value_type as its underlying trait expression rather than
  // the member typedef: clang (through 22) fails to match an out-of-line constrained
  // member of a partial specialization when the requires-clause names a member typedef
  // of the class.
  template<class G = @\exposidnc{error-value-type}@>
    requires(is_copy_constructible_v<remove_cv_t<remove_reference_t<E>>> &&
             is_convertible_v<G, remove_cv_t<remove_reference_t<E>>>)
  constexpr @\exposidnc{error-value-type}@ error_or(G&& def) const&;

  template<class G = @\exposidnc{error-value-type}@>
    requires(is_move_constructible_v<remove_cv_t<remove_reference_t<E>>> &&
             is_convertible_v<G, remove_cv_t<remove_reference_t<E>>>)
  constexpr @\exposidnc{error-value-type}@ error_or(G&& def) &&;

  // -------------------------------------------------------------------------
  // Monadic operations
  // -------------------------------------------------------------------------

  template<class F>
    requires is_constructible_v<E, E&>
  constexpr auto and_then(F&& f) &;
  template<class F>
    requires is_constructible_v<E, E&&>
  constexpr auto and_then(F&& f) &&;
  template<class F>
    requires is_constructible_v<E, const E&>
  constexpr auto and_then(F&& f) const&;
  template<class F>
    requires is_constructible_v<E, const E&&>
  constexpr auto and_then(F&& f) const&&;

  template<class F> constexpr auto or_else(F&& f) &;
  template<class F> constexpr auto or_else(F&& f) &&;
  template<class F> constexpr auto or_else(F&& f) const&;
  template<class F> constexpr auto or_else(F&& f) const&&;

  // transform: f receives T& (value); error propagates as E; result is expected<U, E>
  template<class F>
    requires is_constructible_v<E, E&>
  constexpr auto transform(F&& f) &;
  template<class F>
    requires is_constructible_v<E, E&&>
  constexpr auto transform(F&& f) &&;
  template<class F>
    requires is_constructible_v<E, const E&>
  constexpr auto transform(F&& f) const&;
  template<class F>
    requires is_constructible_v<E, const E&&>
  constexpr auto transform(F&& f) const&&;

  // transform_error: f receives E; value propagates as T&; result is expected<T&, G>
  template<class F> constexpr auto transform_error(F&& f) &;
  template<class F> constexpr auto transform_error(F&& f) &&;
  template<class F> constexpr auto transform_error(F&& f) const&;
  template<class F> constexpr auto transform_error(F&& f) const&&;

  // -------------------------------------------------------------------------
  // Equality operators (hidden friends)
  // -------------------------------------------------------------------------

  template<class T2, class E2>
    requires(!is_void_v<T2>)
  friend constexpr bool operator==(const expected& x, const expected<T2, E2>& y);

  template<class T2>
    requires(!@\exposidnc{is-expected-specialization}@<T2>::value)
  friend constexpr bool operator==(const expected& x, const T2& val);

  template<class E2>
  friend constexpr bool operator==(const expected& x, const unexpected<E2>& e);

private:
  bool @\exposidnc{has-val}@; // exposition only
  union {
    T* @\exposidnc{val}@;             // exposition only
    unexpected<E> @\exposidnc{unex}@; // exposition only
  };
};
\end{codeblock}

\begin{itemdescr}
\pnum
\mandates
A program that instantiates the definition of \tcode{expected<T&, E>} with an \tcode{E} that is not a valid template argument for \tcode{unexpected} is ill-formed. \tcode{T} shall be an object type that is not an array type.

\pnum
\remarks
An object of type \tcode{expected<T&, E>} either represents a reference to an object of type \tcode{T}, or holds an error. Member \tcode{has_val} indicates whether the object represents a reference. When it represents a reference, member \tcode{val} points to the referenced object, which is not owned by the \tcode{expected} object. Otherwise, the error is \tcode{unex.error()}.
\end{itemdescr}

\rSec3[expected.ref.cons]{Constructors}

\indexlibraryctor{expected}%
\begin{itemdecl}
;
\end{itemdecl}

\begin{itemdescr}
\pnum
\remarks
\tcode{expected<T&, E>} has no default constructor: a reference cannot be null, so there is no empty state to default-construct into.
\end{itemdescr}

\indexlibraryctor{expected}%
\begin{itemdecl}
constexpr expected(const expected&) = default;
\end{itemdecl}

\begin{itemdescr}
\pnum
\effects
If \tcode{rhs.has_value()} is \tcode{true}, initializes \tcode{val} with \tcode{rhs.val}, so that \tcode{*this} and \tcode{rhs} refer to the same object; otherwise, initializes \tcode{unex} with \tcode{rhs.unex}.

\pnum
\ensures
\tcode{rhs.has_value() == this->has_value()}.

\pnum
\remarks
This constructor is trivial.
\end{itemdescr}

\indexlibraryctor{expected}%
\begin{itemdecl}
constexpr expected(const expected& rhs) noexcept(is_nothrow_copy_constructible_v<E>);
constexpr expected(expected&& rhs) noexcept(is_nothrow_move_constructible_v<E>);
\end{itemdecl}

\begin{itemdescr}
\pnum
\constraints
\tcode{is_copy_constructible_v<E>} is \tcode{true} and \tcode{is_trivially_copy_constructible_v<E>} is \tcode{false}.

\pnum
\effects
If \tcode{rhs.has_value()} is \tcode{true}, initializes \tcode{val} with \tcode{rhs.val}, so that \tcode{*this} and \tcode{rhs} refer to the same object; otherwise, initializes \tcode{unex} with \tcode{rhs.unex}.

\pnum
\ensures
\tcode{rhs.has_value() == this->has_value()}.

\pnum
\remarks
This constructor is trivial if the corresponding constructor of \tcode{E} is trivial, and is defined as deleted unless \tcode{is_copy_constructible_v<E>} is \tcode{true}.
\end{itemdescr}

\indexlibraryctor{expected}%
\begin{itemdecl}
constexpr expected(expected&&) = default;
\end{itemdecl}

\begin{itemdescr}
\pnum
\effects
If \tcode{rhs.has_value()} is \tcode{true}, initializes \tcode{val} with \tcode{rhs.val}, so that \tcode{*this} and \tcode{rhs} refer to the same object; otherwise, initializes \tcode{unex} with \tcode{std::move(rhs.unex)}.

\pnum
\ensures
\tcode{rhs.has_value() == this->has_value()}.

\pnum
\remarks
This constructor is trivial.
\end{itemdescr}

\indexlibraryctor{expected}%
\begin{itemdecl}
template<class... Args> constexpr expected(in_place_t, Args&&...);
\end{itemdecl}

\begin{itemdescr}
\pnum
\remarks
\tcode{expected<T&, E>} has no in-place value constructor: \tcode{T&} cannot be constructed in-place. Pass a \tcode{U} convertible to \tcode{T&} instead.
\end{itemdescr}

\indexlibraryctor{expected}%
\begin{itemdecl}
template<class U = T>
  requires(!is_same_v<remove_cvref_t<U>, in_place_t> &&
           !is_same_v<remove_cvref_t<U>, expected<T&, E>> &&
           !@\exposidnc{is-unexpected-specialization}@<remove_cvref_t<U>>::value &&
           is_constructible_v<T&, U> && !reference_constructs_from_temporary_v<T&, U>)
constexpr explicit(!is_convertible_v<U, T&>)
    expected(U&& u) noexcept(is_nothrow_constructible_v<T&, U>);
\end{itemdecl}

\begin{itemdescr}
\pnum
\constraints
\tcode{remove_cvref_t<U>} is not \tcode{in_place_t}, \tcode{expected}, or a specialization of \tcode{unexpected}; \tcode{is_constructible_v<T&, U>} is \tcode{true}; and \tcode{reference_constructs_from_temporary_v<T&, U>} is \tcode{false}.

\pnum
\effects
Let \tcode{r} be the lvalue result of \tcode{T& r = std::forward<U>(u);}. Initializes \tcode{val} with \tcode{addressof(r)}.

\pnum
\ensures
\tcode{has_value()} is \tcode{true}.
\end{itemdescr}

\indexlibraryctor{expected}%
\begin{itemdecl}
template<class U>
  requires(reference_constructs_from_temporary_v<T&, U>)
constexpr expected(U&&);
\end{itemdecl}

\begin{itemdescr}
\pnum
\remarks
A constructor for which \tcode{reference_constructs_from_temporary_v<T&, U>} is \tcode{true} — one that would bind \tcode{T&} to a temporary — is defined as deleted.
\end{itemdescr}

\indexlibraryctor{expected}%
\begin{itemdecl}
template<class U, class G>
  requires(!is_reference_v<E> && is_constructible_v<T&, U&> &&
           is_constructible_v<E, const G&> &&
           !reference_constructs_from_temporary_v<T&, U&>)
constexpr explicit(!is_convertible_v<U&, T&> || !is_convertible_v<const G&, E>)
    expected(const expected<U&, G>& rhs);
template<class U, class G>
  requires(!is_reference_v<E> && is_constructible_v<T&, U&> &&
           is_constructible_v<E, G> && !reference_constructs_from_temporary_v<T&, U&>)
constexpr explicit(!is_convertible_v<U&, T&> || !is_convertible_v<G, E>)
    expected(expected<U&, G>&& rhs);
\end{itemdecl}

\begin{itemdescr}
\pnum
\constraints
\tcode{is_constructible_v<T&, U&>} is \tcode{true}; \tcode{reference_constructs_from_temporary_v<T&, U&>} is \tcode{false}; and \tcode{is_constructible_v<E, const G&>} is \tcode{true}.

\pnum
\effects
If \tcode{rhs.has_value()} is \tcode{true}, initializes \tcode{val} with \tcode{addressof(*rhs)}, so that \tcode{*this} refers to the object referred to by \tcode{rhs}; otherwise, initializes \tcode{unex} with the error of \tcode{rhs}. No object referred to by \tcode{rhs} is moved from.
\end{itemdescr}

\indexlibraryctor{expected}%
\begin{itemdecl}
template<class U, class G>
  requires(is_reference_v<E> && is_reference_v<G> && is_constructible_v<T&, U&> &&
           is_convertible_v<G, E> && !reference_constructs_from_temporary_v<T&, U&>)
constexpr explicit(!is_convertible_v<U&, T&> || !is_convertible_v<G, E>)
    expected(const expected<U&, G>& rhs);
template<class U, class G>
  requires(is_reference_v<E> && is_reference_v<G> && is_constructible_v<T&, U&> &&
           is_convertible_v<G, E> && !reference_constructs_from_temporary_v<T&, U&>)
constexpr explicit(!is_convertible_v<U&, T&> || !is_convertible_v<G, E>)
    expected(expected<U&, G>&& rhs);
\end{itemdecl}

\begin{itemdescr}
\pnum
\constraints
\tcode{is_constructible_v<T&, U&>} is \tcode{true}; \tcode{reference_constructs_from_temporary_v<T&, U&>} is \tcode{false}; \tcode{G} is a reference type; and \tcode{is_convertible_v<G, E>} is \tcode{true}.

\pnum
\effects
If \tcode{rhs.has_value()} is \tcode{true}, initializes \tcode{val} with \tcode{addressof(*rhs)}, so that \tcode{*this} refers to the object referred to by \tcode{rhs}; otherwise, initializes \tcode{unex} with the error of \tcode{rhs}. No object referred to by \tcode{rhs} is moved from.
\end{itemdescr}

\indexlibraryctor{expected}%
\begin{itemdecl}
template<class G>
  requires(!is_reference_v<E> && is_constructible_v<E, const G&>)
constexpr explicit(!is_convertible_v<const G&, E>) expected(const unexpected<G>& e);
template<class G>
  requires(!is_reference_v<E> && is_constructible_v<E, G>)
constexpr explicit(!is_convertible_v<G, E>) expected(unexpected<G>&& e);
\end{itemdecl}

\begin{itemdescr}
\pnum
\constraints
\tcode{is_constructible_v<E, const G&>} is \tcode{true}.

\pnum
\effects
Initializes \tcode{unex} with the error of \tcode{e}.

\pnum
\ensures
\tcode{has_value()} is \tcode{false}.
\end{itemdescr}

\indexlibraryctor{expected}%
\begin{itemdecl}
template<class G>
  requires(is_reference_v<E> && is_reference_v<G> && is_constructible_v<E, G> &&
           !reference_constructs_from_temporary_v<E, G>)
constexpr explicit(!is_convertible_v<G, E>) expected(const unexpected<G>& e) noexcept;
template<class G>
  requires(is_reference_v<E> && is_reference_v<G> && is_constructible_v<E, G> &&
           !reference_constructs_from_temporary_v<E, G>)
constexpr explicit(!is_convertible_v<G, E>) expected(unexpected<G>&& e) noexcept;
\end{itemdecl}

\begin{itemdescr}
\pnum
\constraints
\tcode{is_reference_v<G>} is \tcode{true}; \tcode{is_convertible_v<G, E>} is \tcode{true}; and \tcode{reference_constructs_from_temporary_v<E, G>} is \tcode{false}.

\pnum
\effects
Initializes \tcode{unex} with the error of \tcode{e}.

\pnum
\ensures
\tcode{has_value()} is \tcode{false}.

\pnum
\remarks
This constructor never throws: the referent is bound, not copied.
\end{itemdescr}

\indexlibraryctor{expected}%
\begin{itemdecl}
template<class G>
  requires(is_reference_v<E> && !is_reference_v<G>)
constexpr expected(const unexpected<G>&);
\end{itemdecl}

\begin{itemdescr}
\pnum
\remarks
When \tcode{E} is a reference type, an overload taking \tcode{unexpected<G>} for a non-reference \tcode{G} is defined as deleted: the referent would live inside the (possibly temporary) source \tcode{unexpected<G>} object, and binding \tcode{E&} to it would dangle. Use an \tcode{unexpected<E&>} holding an external object instead.
\end{itemdescr}

\indexlibraryctor{expected}%
\begin{itemdecl}
template<class... Args>
  requires(is_constructible_v<E, Args...> && !@\exposidnc{unexpect-dangles-v}@<E, Args...>)
constexpr explicit expected(unexpect_t, Args&&... args);
\end{itemdecl}

\begin{itemdescr}
\pnum
\constraints
\tcode{is_constructible_v<E, Args...>} is \tcode{true}.

\pnum
\effects
Direct-non-list-initializes \tcode{unex} with \tcode{in_place} and \tcode{std::forward<Args>(args)...}.

\pnum
\ensures
\tcode{has_value()} is \tcode{false}.
\end{itemdescr}

\indexlibraryctor{expected}%
\begin{itemdecl}
template<class... Args>
  requires(@\exposidnc{unexpect-dangles-v}@<E, Args...>)
constexpr expected(unexpect_t, Args&&...);
\end{itemdecl}

\begin{itemdescr}
\pnum
\remarks
When \tcode{E} is a reference type, an overload with the same parameter types is defined as deleted if the single argument would bind \tcode{E&} to a temporary, or if it is otherwise not usable to construct \tcode{E}.
\end{itemdescr}

\indexlibraryctor{expected}%
\begin{itemdecl}
template<class U, class... Args>
  requires(!is_reference_v<E> && is_constructible_v<E, initializer_list<U>&, Args...>)
constexpr explicit expected(unexpect_t, initializer_list<U> il, Args&&... args);
\end{itemdecl}

\begin{itemdescr}
\pnum
\constraints
\tcode{is_reference_v<E>} is \tcode{false}, and \tcode{is_constructible_v<E, initializer_list<U>&, Args...>} is \tcode{true}.

\pnum
\effects
Direct-non-list-initializes \tcode{unex} with \tcode{in_place}, \tcode{il}, and \tcode{std::forward<Args>(args)...}.

\pnum
\ensures
\tcode{has_value()} is \tcode{false}.

\pnum
\remarks
An overload with the same parameter types is defined as deleted when \tcode{E} is an lvalue reference type. An initializer list cannot provide the required long-lived error referent.
\end{itemdescr}

\indexlibraryctor{expected}%
\begin{itemdecl}
template<class U, class... Args>
  requires is_reference_v<E>
constexpr expected(unexpect_t, initializer_list<U>, Args&&...);
\end{itemdecl}

\begin{itemdescr}
\pnum
\remarks
An overload with the same parameter types is defined as deleted when \tcode{E} is an lvalue reference type. An initializer list cannot provide the required long-lived error referent.
\end{itemdescr}

\rSec3[expected.ref.dtor]{Destructor}

\indexlibrarydtor{expected}%
\begin{itemdecl}
constexpr ~expected();
\end{itemdecl}

\begin{itemdescr}
\pnum
\constraints
\tcode{is_trivially_destructible_v<E>} is \tcode{true}.

\pnum
\effects
None: \tcode{*this} never owns the object it refers to; \tcode{T} is never destroyed.

\pnum
\remarks
This destructor is trivial.
\end{itemdescr}

\indexlibrarydtor{expected}%
\begin{itemdecl}
constexpr ~expected();
\end{itemdecl}

\begin{itemdescr}
\pnum
\constraints
\tcode{is_trivially_destructible_v<E>} is \tcode{false}.

\pnum
\effects
If \tcode{has_value()} is \tcode{false}, destroys \tcode{unex}. \tcode{T} is not destroyed; \tcode{*this} never owns the object it refers to.

\pnum
\remarks
This destructor is trivial if \tcode{E} is trivially destructible.
\end{itemdescr}

\rSec3[expected.ref.assign]{Assignment}

\indexlibrarymember{operator=}{expected}%
\begin{itemdecl}
constexpr expected& operator=(const expected&);
\end{itemdecl}

\begin{itemdescr}
\pnum
\constraints
\tcode{is_trivially_copy_constructible_v<E>} is \tcode{true}, \tcode{is_trivially_copy_assignable_v<E>} is \tcode{true}, and \tcode{is_trivially_destructible_v<E>} is \tcode{true}.

\pnum
\effects
If \tcode{rhs.has_value()} is \tcode{true}: if \tcode{has_value()} is \tcode{true}, assigns \tcode{rhs.val} to \tcode{val}; otherwise destroys \tcode{unex} and initializes \tcode{val} with \tcode{rhs.val}. If \tcode{rhs.has_value()} is \tcode{false}, the error of \tcode{rhs} is assigned to or used to initialize \tcode{unex}, as for the primary template. In every case \tcode{*this} comes to refer to the object \tcode{rhs} refers to, or to hold the error of \tcode{rhs}.

\pnum
\returns
\tcode{*this}.

\pnum
\remarks
Assignment rebinds: assigning to an \tcode{expected<T&, E>} that holds a value changes which object it refers to. It never assigns through to the referent. This operator is trivial.
\end{itemdescr}

\indexlibrarymember{operator=}{expected}%
\begin{itemdecl}
constexpr expected& operator=(const expected& rhs) noexcept(
    is_nothrow_copy_constructible_v<E> && is_nothrow_copy_assignable_v<E>);
constexpr expected& operator=(expected&& rhs) noexcept(
    is_nothrow_move_constructible_v<E> && is_nothrow_move_assignable_v<E>);
\end{itemdecl}

\begin{itemdescr}
\pnum
\constraints
\tcode{is_reference_v<E> || (is_copy_constructible_v<E> && is_copy_assignable_v<E>)} is \tcode{true} and \tcode{(is_trivially_copy_constructible_v<E> && is_trivially_copy_assignable_v<E> &&
                           is_trivially_destructible_v<E>)} is \tcode{false}.

\pnum
\effects
If \tcode{rhs.has_value()} is \tcode{true}: if \tcode{has_value()} is \tcode{true}, assigns \tcode{rhs.val} to \tcode{val}; otherwise destroys \tcode{unex} and initializes \tcode{val} with \tcode{rhs.val}. If \tcode{rhs.has_value()} is \tcode{false}, the error of \tcode{rhs} is assigned to or used to initialize \tcode{unex}, as for the primary template. In every case \tcode{*this} comes to refer to the object \tcode{rhs} refers to, or to hold the error of \tcode{rhs}.

\pnum
\returns
\tcode{*this}.

\pnum
\remarks
This operator is defined as deleted unless \tcode{is_copy_assignable_v<E>} is \tcode{true} and \tcode{is_copy_constructible_v<E>} is \tcode{true}.
\end{itemdescr}

\indexlibrarymember{operator=}{expected}%
\begin{itemdecl}
constexpr expected& operator=(expected&&) noexcept;
\end{itemdecl}

\begin{itemdescr}
\pnum
\constraints
\tcode{is_trivially_move_constructible_v<E>} is \tcode{true}, \tcode{is_trivially_move_assignable_v<E>} is \tcode{true}, and \tcode{is_trivially_destructible_v<E>} is \tcode{true}.

\pnum
\effects
If \tcode{rhs.has_value()} is \tcode{true}: if \tcode{has_value()} is \tcode{true}, assigns \tcode{rhs.val} to \tcode{val}; otherwise destroys \tcode{unex} and initializes \tcode{val} with \tcode{rhs.val}. If \tcode{rhs.has_value()} is \tcode{false}, the error of \tcode{rhs} is assigned to or used to initialize \tcode{unex}, as for the primary template.

\pnum
\returns
\tcode{*this}.

\pnum
\remarks
This operator is trivial.
\end{itemdescr}

\indexlibrarymember{operator=}{expected}%
\begin{itemdecl}
template<class U = T>
  requires(!is_same_v<remove_cvref_t<U>, expected<T&, E>> &&
           !@\exposidnc{is-unexpected-specialization}@<remove_cvref_t<U>>::value &&
           is_constructible_v<T&, U> && !reference_constructs_from_temporary_v<T&, U>)
constexpr expected& operator=(U&& u);
\end{itemdecl}

\begin{itemdescr}
\pnum
\constraints
\tcode{remove_cvref_t<U>} is neither \tcode{expected} nor a specialization of \tcode{unexpected}, \tcode{is_constructible_v<T&, U>} is \tcode{true}, and \tcode{reference_constructs_from_temporary_v<T&, U>} is \tcode{false}.

\pnum
\effects
Let \tcode{r} be the lvalue result of \tcode{T& r = std::forward<U>(u);}. If \tcode{has_value()} is \tcode{true}, assigns \tcode{addressof(r)} to \tcode{val}. Otherwise, destroys \tcode{unex}, initializes \tcode{val} with \tcode{addressof(r)}, and sets \tcode{has_val} to \tcode{true}; if binding \tcode{r} throws, \tcode{*this} is left unchanged.

\pnum
\returns
\tcode{*this}.
\end{itemdescr}

\indexlibrarymember{operator=}{expected}%
\begin{itemdecl}
template<class G>
  requires(!is_reference_v<E> && is_constructible_v<E, const G&> &&
           is_assignable_v<E&, const G&>)
constexpr expected& operator=(const unexpected<G>& e);
template<class G>
  requires(!is_reference_v<E> && is_constructible_v<E, G> && is_assignable_v<E&, G>)
constexpr expected& operator=(unexpected<G>&& e);
\end{itemdecl}

\begin{itemdescr}
\pnum
\constraints
\tcode{is_constructible_v<E, const G&>} is \tcode{true} and \tcode{is_assignable_v<E&, const G&>} is \tcode{true}.

\pnum
\effects
Makes \tcode{*this} hold the error of \tcode{e}, reinitializing \tcode{unex} from \tcode{e} rather than assigning through it.

\pnum
\returns
\tcode{*this}.
\end{itemdescr}

\indexlibrarymember{operator=}{expected}%
\begin{itemdecl}
template<class G>
  requires(is_reference_v<E> && is_reference_v<G> && is_constructible_v<E, G> &&
           !reference_constructs_from_temporary_v<E, G>)
constexpr expected& operator=(const unexpected<G>& e);
template<class G>
  requires(is_reference_v<E> && is_reference_v<G> && is_constructible_v<E, G> &&
           !reference_constructs_from_temporary_v<E, G>)
constexpr expected& operator=(unexpected<G>&& e);
\end{itemdecl}

\begin{itemdescr}
\pnum
\constraints
\tcode{is_reference_v<G>} is \tcode{true}; \tcode{is_convertible_v<G, E>} is \tcode{true}; and \tcode{reference_constructs_from_temporary_v<E, G>} is \tcode{false}.

\pnum
\effects
Makes \tcode{*this} hold the error of \tcode{e}, reinitializing \tcode{unex} from \tcode{e} rather than assigning through it. \tcode{unex.error()} thereafter refers to the same object as \tcode{e.error()}; the previously referenced object, if any, is not modified.

\pnum
\returns
\tcode{*this}.

\pnum
\remarks
This operator never throws: the referent is bound, not copied.
\end{itemdescr}

\indexlibrarymember{operator=}{expected}%
\begin{itemdecl}
template<class G>
  requires(is_reference_v<E> && !is_reference_v<G>)
constexpr expected& operator=(const unexpected<G>&);
\end{itemdecl}

\begin{itemdescr}
\pnum
\remarks
When \tcode{E} is a reference type, an overload taking \tcode{unexpected<G>} for a non-reference \tcode{G} is defined as deleted: it would rebind \tcode{unex} to \tcode{unexpected<G>}'s temporary storage.
\end{itemdescr}

\indexlibrarymember{emplace}{expected}%
\begin{itemdecl}
template<class U = T>
  requires(is_constructible_v<T&, U> && !reference_constructs_from_temporary_v<T&, U>)
constexpr T& emplace(U&& u) noexcept(is_nothrow_constructible_v<T&, U>);
\end{itemdecl}

\begin{itemdescr}
\pnum
\constraints
\tcode{is_constructible_v<T&, U>} is \tcode{true} and \tcode{reference_constructs_from_temporary_v<T&, U>} is \tcode{false}.

\pnum
\effects
Rebinds \tcode{*this} to refer to the object bound by \tcode{T& r = std::forward<U>(u);}: if \tcode{has_value()} is \tcode{false}, destroys \tcode{unex} first. Sets \tcode{val} to \tcode{addressof(r)} and \tcode{has_val} to \tcode{true}.

\pnum
\returns
\tcode{*val}.
\end{itemdescr}

\rSec3[expected.ref.swap]{Swap}

\indexlibrarymember{swap}{expected}%
\begin{itemdecl}
constexpr void swap(expected& rhs) noexcept(is_nothrow_move_constructible_v<E> &&
                                            (is_reference_v<E> ||
                                             is_nothrow_swappable_v<E>));
\end{itemdecl}

\begin{itemdescr}
\pnum
\constraints
\tcode{is_reference_v<E> || is_swappable_v<E>} is \tcode{true} and \tcode{is_move_constructible_v<E>} is \tcode{true}.

\pnum
\effects
Exchanges the states of \tcode{*this} and \tcode{rhs}. When both hold values, exchanges \tcode{val} and \tcode{rhs.val} — the referenced objects are not swapped. Otherwise behaves as the primary template's \tcode{swap} does for the error.

\pnum
\remarks
The exception specification is equivalent to \tcode{is_nothrow_move_constructible_v<E> && (is_reference_v<E> || is_nothrow_swappable_v<E>)}.
\end{itemdescr}

\indexlibraryglobal{swap}%
\begin{itemdecl}
friend constexpr void swap(expected& x, expected& y) noexcept(noexcept(x.swap(y)));
\end{itemdecl}

\begin{itemdescr}
\pnum
\constraints
\tcode{is_reference_v<E> || is_swappable_v<E>} is \tcode{true} and \tcode{is_move_constructible_v<E>} is \tcode{true}.

\pnum
\effects
Equivalent to \tcode{x.swap(y)}.
\end{itemdescr}

\rSec3[expected.ref.obs]{Observers}

\indexlibrarymember{operator->}{expected}%
\begin{itemdecl}
constexpr T* operator->() const noexcept;
\end{itemdecl}

\begin{itemdescr}
\pnum
\expects
\tcode{has_value()} is \tcode{true}.

\pnum
\returns
\tcode{val}.

\pnum
\remarks
This is a \tcode{const} member function that returns a non-\tcode{const} \tcode{T*}; the constness of \tcode{*this} does not propagate to the referenced object. For deep \tcode{const}, use \tcode{expected<const T&, E>}.
\end{itemdescr}

\indexlibrarymember{operator*}{expected}%
\begin{itemdecl}
constexpr T& operator*() const noexcept;
\end{itemdecl}

\begin{itemdescr}
\pnum
\expects
\tcode{has_value()} is \tcode{true}.

\pnum
\returns
\tcode{*val}.

\pnum
\remarks
This is a \tcode{const} member function that returns a non-\tcode{const} \tcode{T&}; the constness of \tcode{*this} does not propagate to the referenced object. For deep \tcode{const}, use \tcode{expected<const T&, E>}.
\end{itemdescr}

\indexlibrarymember{operator bool}{expected}%
\indexlibrarymember{has_value}{expected}%
\begin{itemdecl}
constexpr explicit operator bool() const noexcept;
constexpr bool has_value() const noexcept;
\end{itemdecl}

\begin{itemdescr}
\pnum
\returns
\tcode{has_val}.
\end{itemdescr}

\indexlibrarymember{value}{expected}%
\begin{itemdecl}
constexpr T& value() const&;
\end{itemdecl}

\begin{itemdescr}
\pnum
\mandates
\tcode{is_copy_constructible_v<\exposid{error-value-type}>} is \tcode{true}.

\pnum
\returns
\tcode{*val} if \tcode{has_value()} is \tcode{true}.

\pnum
\throws
\tcode{bad_expected_access(as_const(error()))} if \tcode{has_value()} is \tcode{false}.
\end{itemdescr}

\indexlibrarymember{value}{expected}%
\begin{itemdecl}
constexpr T& value() &&;
\end{itemdecl}

\begin{itemdescr}
\pnum
\returns
\tcode{*val} if \tcode{has_value()} is \tcode{true}.

\pnum
\throws
\tcode{bad_expected_access(std::move(error()))} if \tcode{has_value()} is \tcode{false}.
\end{itemdescr}

\indexlibrarymember{error}{expected}%
\begin{itemdecl}
constexpr const E& error() const& noexcept;
constexpr E& error() & noexcept;
\end{itemdecl}

\begin{itemdescr}
\pnum
\expects
\tcode{has_value()} is \tcode{false}.

\pnum
\returns
\tcode{unex.error()}.
\end{itemdescr}

\indexlibrarymember{error}{expected}%
\begin{itemdecl}
constexpr const E&& error() const&& noexcept;
constexpr E&& error() && noexcept;
\end{itemdecl}

\begin{itemdescr}
\pnum
\expects
\tcode{has_value()} is \tcode{false}.

\pnum
\returns
\tcode{std::move(unex).error()}.
\end{itemdescr}

\indexlibrarymember{value_or}{expected}%
\begin{itemdecl}
template<class U = remove_cv_t<T>>
  requires(is_object_v<T> && !is_array_v<T>)
constexpr remove_cv_t<T> value_or(U&& def) const;
\end{itemdecl}

\begin{itemdescr}
\pnum
\mandates
\tcode{is_convertible_v<T&, remove_cv_t<T>>} and \tcode{is_convertible_v<U, remove_cv_t<T>>} are \tcode{true}.

\pnum
\returns
\tcode{has_value() ? static_cast<remove_cv_t<T>>(*val) : static_cast<remove_cv_t<T>>(std::forward<U>(def))}. The result is an object, never a reference.
\end{itemdescr}

\indexlibrarymember{error_or}{expected}%
\begin{itemdecl}
template<class G = @\exposidnc{error-value-type}@>
  requires(is_copy_constructible_v<remove_cv_t<remove_reference_t<E>>> &&
           is_convertible_v<G, remove_cv_t<remove_reference_t<E>>>)
constexpr @\exposidnc{error-value-type}@ error_or(G&& def) const&;
\end{itemdecl}

\begin{itemdescr}
\pnum
\returns
\tcode{std::forward<G>(def)} if \tcode{has_value()} is \tcode{true}, \tcode{error()} otherwise. The result is an object, never a reference.
\end{itemdescr}

\indexlibrarymember{error_or}{expected}%
\begin{itemdecl}
template<class G = @\exposidnc{error-value-type}@>
  requires(is_move_constructible_v<remove_cv_t<remove_reference_t<E>>> &&
           is_convertible_v<G, remove_cv_t<remove_reference_t<E>>>)
constexpr @\exposidnc{error-value-type}@ error_or(G&& def) &&;
\end{itemdecl}

\begin{itemdescr}
\pnum
\returns
\tcode{std::forward<G>(def)} if \tcode{has_value()} is \tcode{true}, \tcode{std::move(error())} otherwise. The result is an object, never a reference.
\end{itemdescr}

\rSec3[expected.ref.monadic]{Monadic operations}

\indexlibrarymember{and_then}{expected}%
\begin{itemdecl}
template<class F>
  requires is_constructible_v<E, E&>
constexpr auto and_then(F&& f) &;
template<class F>
  requires is_constructible_v<E, const E&>
constexpr auto and_then(F&& f) const&;
\end{itemdecl}

\begin{itemdescr}
\pnum
\constraints
\tcode{is_constructible_v<E, decltype(error())>} is \tcode{true}.

\pnum
\mandates
\tcode{remove_cvref_t<invoke_result_t<F, T&>>} is a specialization of \tcode{expected} and its \tcode{error_type} is the same type as \tcode{E}.

\pnum
\effects
Equivalent to: \tcode{if (has_value()) return invoke(std::forward<F>(f), *val); else return U(unexpect, error());} where \tcode{U} is \tcode{remove_cvref_t<invoke_result_t<F, T&>>}.

\pnum
\remarks
The member templates \tcode{and_then}, \tcode{or_else}, \tcode{transform}, and \tcode{transform_error} behave as specified for the primary template, with one difference: the value is passed to the callable as \tcode{T&} for every ref-qualification of \tcode{*this}. An rvalue \tcode{expected<T&, E>} does not pass its referent as an rvalue; the object referred to is never moved from by these operations.
\end{itemdescr}

\indexlibrarymember{and_then}{expected}%
\begin{itemdecl}
template<class F>
  requires is_constructible_v<E, E&&>
constexpr auto and_then(F&& f) &&;
template<class F>
  requires is_constructible_v<E, const E&&>
constexpr auto and_then(F&& f) const&&;
\end{itemdecl}

\begin{itemdescr}
\pnum
\constraints
\tcode{is_constructible_v<E, decltype(std::move(error()))>} is \tcode{true}.

\pnum
\mandates
\tcode{remove_cvref_t<invoke_result_t<F, T&>>} is a specialization of \tcode{expected} and its \tcode{error_type} is the same type as \tcode{E}.

\pnum
\effects
Equivalent to: \tcode{if (has_value()) return invoke(std::forward<F>(f), *val); else return U(unexpect, std::move(error()));} where \tcode{U} is \tcode{remove_cvref_t<invoke_result_t<F, T&>>}.
\end{itemdescr}

\indexlibrarymember{or_else}{expected}%
\begin{itemdecl}
template<class F> constexpr auto or_else(F&& f) &;
template<class F> constexpr auto or_else(F&& f) const&;
\end{itemdecl}

\begin{itemdescr}
\pnum
\mandates
\tcode{remove_cvref_t<invoke_result_t<F, decltype(error())>>} is a specialization of \tcode{expected} and its \tcode{value_type} is the same type as \tcode{T&}.

\pnum
\effects
Equivalent to: \tcode{if (has_value()) return G(*val); else return invoke(std::forward<F>(f), error());} where \tcode{G} is \tcode{remove_cvref_t<invoke_result_t<F, decltype(error())>>}.
\end{itemdescr}

\indexlibrarymember{or_else}{expected}%
\begin{itemdecl}
template<class F> constexpr auto or_else(F&& f) &&;
template<class F> constexpr auto or_else(F&& f) const&&;
\end{itemdecl}

\begin{itemdescr}
\pnum
\mandates
\tcode{remove_cvref_t<invoke_result_t<F, decltype(std::move(error()))>>} is a specialization of \tcode{expected} and its \tcode{value_type} is the same type as \tcode{T&}.

\pnum
\effects
Equivalent to: \tcode{if (has_value()) return G(*val); else return invoke(std::forward<F>(f), std::move(error()));} where \tcode{G} is \tcode{remove_cvref_t<invoke_result_t<F, decltype(std::move(error()))>>}.
\end{itemdescr}

\indexlibrarymember{transform}{expected}%
\begin{itemdecl}
template<class F>
  requires is_constructible_v<E, E&>
constexpr auto transform(F&& f) &;
template<class F>
  requires is_constructible_v<E, const E&>
constexpr auto transform(F&& f) const&;
\end{itemdecl}

\begin{itemdescr}
\pnum
\constraints
\tcode{is_constructible_v<E, decltype(error())>} is \tcode{true}.

\pnum
\effects
Equivalent to: \tcode{if (!has_value()) return U(unexpect, error()); else return expected<U2, E>(in_place, invoke(std::forward<F>(f), *val));} where \tcode{U2} is \tcode{remove_cv_t<invoke_result_t<F, T&>>} and \tcode{U} is \tcode{expected<U2, E>}.
\end{itemdescr}

\indexlibrarymember{transform}{expected}%
\begin{itemdecl}
template<class F>
  requires is_constructible_v<E, E&&>
constexpr auto transform(F&& f) &&;
template<class F>
  requires is_constructible_v<E, const E&&>
constexpr auto transform(F&& f) const&&;
\end{itemdecl}

\begin{itemdescr}
\pnum
\constraints
\tcode{is_constructible_v<E, decltype(std::move(error()))>} is \tcode{true}.

\pnum
\effects
Equivalent to: \tcode{if (!has_value()) return U(unexpect, std::move(error())); else return expected<U2, E>(in_place, invoke(std::forward<F>(f), *val));} where \tcode{U2} is \tcode{remove_cv_t<invoke_result_t<F, T&>>} and \tcode{U} is \tcode{expected<U2, E>}.
\end{itemdescr}

\indexlibrarymember{transform_error}{expected}%
\begin{itemdecl}
template<class F> constexpr auto transform_error(F&& f) &;
template<class F> constexpr auto transform_error(F&& f) const&;
\end{itemdecl}

\begin{itemdescr}
\pnum
\mandates
\begin{itemize}
\item \tcode{is_object_v<G>} is \tcode{true},
\item \tcode{is_array_v<G>} is \tcode{false},
\item \tcode{is_same_v<G, remove_cv_t<G>>} is \tcode{true},
\item \tcode{\exposid{is-unexpected-specialization}<G>::value} is \tcode{false}.
\end{itemize}

\pnum
\effects
Equivalent to: \tcode{if (has_value()) return G(*val); else return expected<T&, G2>(unexpect, invoke(std::forward<F>(f), error()));} where \tcode{G2} is \tcode{remove_cv_t<invoke_result_t<F, decltype(error())>>} and \tcode{G} is \tcode{expected<T&, G2>}.
\end{itemdescr}

\indexlibrarymember{transform_error}{expected}%
\begin{itemdecl}
template<class F> constexpr auto transform_error(F&& f) &&;
template<class F> constexpr auto transform_error(F&& f) const&&;
\end{itemdecl}

\begin{itemdescr}
\pnum
\mandates
\begin{itemize}
\item \tcode{is_object_v<G>} is \tcode{true},
\item \tcode{is_array_v<G>} is \tcode{false},
\item \tcode{is_same_v<G, remove_cv_t<G>>} is \tcode{true},
\item \tcode{\exposid{is-unexpected-specialization}<G>::value} is \tcode{false}.
\end{itemize}

\pnum
\effects
Equivalent to: \tcode{if (has_value()) return G(*val); else return expected<T&, G2>(unexpect, invoke(std::forward<F>(f), std::move(error())));} where \tcode{G2} is \tcode{remove_cv_t<invoke_result_t<F, decltype(std::move(error()))>>} and \tcode{G} is \tcode{expected<T&, G2>}.
\end{itemdescr}

\rSec3[expected.ref.eq]{Equality operators}

\indexlibraryglobal{operator==}%
\begin{itemdecl}
template<class T2, class E2>
  requires(!is_void_v<T2>)
friend constexpr bool operator==(const expected& x, const expected<T2, E2>& y);
\end{itemdecl}

\begin{itemdescr}
\pnum
\mandates
\tcode{!is_void_v<T2>} is \tcode{true}. The expression \tcode{*x == *y} is well-formed and its result is convertible to \tcode{bool}. The expression \tcode{x.error() == y.error()} is well-formed and its result is convertible to \tcode{bool}.

\pnum
\returns
If \tcode{x.has_value() != y.has_value()}, \tcode{false}; otherwise, if \tcode{x.has_value()} is \tcode{true}, \tcode{*x == *y}; otherwise \tcode{x.error() == y.error()}.

\pnum
\remarks
The equality operators behave as specified for the primary template, comparing referents through \tcode{operator*} and errors through \tcode{error()}.
\end{itemdescr}

\indexlibraryglobal{operator==}%
\begin{itemdecl}
template<class T2>
  requires(!@\exposidnc{is-expected-specialization}@<T2>::value)
friend constexpr bool operator==(const expected& x, const T2& val);
\end{itemdecl}

\begin{itemdescr}
\pnum
\mandates
\tcode{T2} is not a specialization of \tcode{expected}. The expression \tcode{*x == val} is well-formed and its result is convertible to \tcode{bool}.

\pnum
\returns
\tcode{x.has_value() && static_cast<bool>(*x == val)}.
\end{itemdescr}

\indexlibraryglobal{operator==}%
\begin{itemdecl}
template<class E2>
friend constexpr bool operator==(const expected& x, const unexpected<E2>& e);
\end{itemdecl}

\begin{itemdescr}
\pnum
\mandates
The expression \tcode{x.error() == e.error()} is well-formed and its result is convertible to \tcode{bool}.

\pnum
\returns
\tcode{!x.has_value() && static_cast<bool>(x.error() == e.error())}.
\end{itemdescr}
